%% APPENDIX E -- radio-frequency interference screening.
%% Owner: appendix-triage.  Carries \label{app:rfi}.
\section{Radio-frequency interference screening}
\label{app:rfi}

Interference is the first alternative explanation for an unattributed
crossing, so the screening this survey has and the coverage it lacks are both
set out here. Millimetre interferometry is an unusually poor environment for
terrestrial interference, though not an empty one: ITU Radio Regulations
Article~5 \citep{ITU2024} carries active allocations across the searched
bands, while an interferometer is not a total-power detector, so a signal not
in the far field at the phase centre acquires a rapidly varying geometric
phase across the array and is suppressed in the correlation rather than
imaged. The documented hazard is spaceborne cloud radar in the
\RadarLo--\RadarHi\,GHz Earth exploration-satellite allocation
\citep{CloudSat2002,EarthCARE2023,ALMAMemo504}: \RadarNWin{} of the
\RadarNTot{} census windows overlap it, carrying \RadarNCross{} threshold
crossings and \RadarNStage{} rank-passing events, so no result here rests on
it. \emph{The five tests of Table~\ref{tab:rfitests} carry none of the signatures
interference leaves, and their coverage is uneven.} Each is computed for every
searched window of the census, as is the \NCtrl-position rank screen, while
recurrence reaches only the crossings with a covering repeat block; and
because the observatory's own correlator flagging was not retained window by
window, \emph{no per-channel interference flag of our own exists, and none can
be recovered} -- which is also why the five populations differ, each test
needing a different quantity to have been recorded. The clustering the
crossings do show is the attribution seen from the other side: the one
concentration a reader will notice, \TlBandAll{} crossings between
\TlBandLo{} and \TlBandHi\,GHz, is the carbon monoxide tuning, \TlBandObs{} of
the \TlBandNSub{} unattributed crossings the census and the reserved hold-out
together place in a searched window against
\TlBandNull$\,\pm\,$\TlBandNullSd{} expected from where the spectral windows
were tuned, and that interval holds \TlBandNLine{} telluric transitions. The
mask must precede the occupancy screen, because \RfiOccMaskN{} of the
\RfiOccSigN{} above-trigger windows sit on a masked transition and a carbon
monoxide line falls at nearly the same sky frequency toward every nearby star,
which is the coincidence that screen is looking for; applied, it reads
\RfiOccMskObs{} against \RfiOccMskNull$\,\pm\,$\RfiOccMskNullSd. The
atmospheric screen is not sharp where it fires -- the one crossing it removes
sits in a window tuned on the line, \ScNearWinPct{} per cent of whose channels
lie within the same radius of the same transition -- so it establishes that
this window cannot separate the two and not that the channel is atmospheric;
that crossing remains among the \MkUnattr{} unattributed, and the same radius
is taken off the control side as one pooled surviving fraction, the larger of
the two, so the residual asymmetry runs against the null of
Appendix~\ref{app:falsealarm} and not for it. The fifth test bounds the one
mechanism that could put a crossing where no line is catalogued: the first
local oscillator is not delivered, but each block's own spectral windows pin
$\nu_{\rm LO1}$ to a median \ImLoWidthBlkGHz\,GHz and so place every
crossing's image \ImGapLoGHz--\ImGapHiGHz\,GHz away, never inside its own
window, while \ImNNoParentWide{} of the \ImNScreen{} survive on frequency
alone with the interval taken \ImWideGHz\,GHz wider at each end, so it is a
bound and not a clearance, and \ImNDsb{} crossings lie where the receivers are
double-sideband and no such bound could be drawn. Where a carbon monoxide
isotopologue does fall in the image band, \ScImNHitCross{} crossings lie at
the image position against \ScImExp{} expected and neither parent attributes
its crossing: each was observed at its own frequency in the same block with
nothing at the trigger, the largest amplitude \ScImUnRatioA{} and
\ScImUnRatioB{} times the crossing's own where a leak at the \ImIrrDb\,dB
image rejection ALMA specifies would need \ImIrrFac{} times as much.

\begin{table}
\centering
\footnotesize
%% ★ v4.26 (r25-captions): 4 printed lines -> the author's 3.  The reason the
%% five populations differ MOVED to the paragraph above, where the record
%% that is or is not retained is the subject.
\caption{The five interference tests, each of which could have failed. $p$ is
against the same peaks redrawn on their own channel grids.}
\label{tab:rfitests}
\begin{tabular}{@{}p{0.20\columnwidth}p{0.22\columnwidth}p{0.33\columnwidth}r@{}}
\hline
Test & Population & Result & $p$ \\
\hline
Clustering in sky frequency &
\AppMNCrossF{} census crossings with a released peak &
\AppMTubeObs{} within $\pm\DispoHalfKms$\,km\,s$^{-1}$ of a searched
transition against \AppMTubeNull{} expected: they cluster on the survey's own
lines &
\AppMTubeP \\
Crowding on one another &
\QcNQP{} unattributed, quality-passing, of the census &
no crowding at any of six separations from \QcTolLo{} to \QcTolHi\,MHz, where
the \QcNAttr{} attributed crossings crowd at every one &
$\geq\QcPQPMin$ \\
Atmospheric transitions &
\TlNCross{} crossings of the census and the hold-out, against \TlNLine{}
telluric lines &
\ScNInRad{} inside the \ScRadMHz\,MHz screening radius in the topocentric
frame, flagged as atmospheric and removed from both sides of the chance
comparison while remaining among the unattributed; \TlTubeObs{} inside
$\pm\TlTubeHalf$\,km\,s$^{-1}$ against \TlTubeNull{} expected &
-- \\
Fixed intermediate frequency &
\IfPopNWin{} census windows with a released peak frequency &
peak positions within their own windows uniform, median \RfiFracMed &
-- \\
Cross-target occupancy &
\RfiOccNWin{} extracted windows toward \RfiOccNStars{} stars, every block set
of Table~\ref{tab:blockledger} together &
\RfiOccObs{} cross-star coincidences within \RfiOccTol\,MHz against
\RfiOccNull$\,\pm\,$\RfiOccNullSd{} expected; \emph{below} the null over the
\RfiOccSigN{} above-trigger windows, \RfiOccSigObs{} against
\RfiOccSigNull$\,\pm\,$\RfiOccSigNullSd &
\RfiOccP \\
Image sideband &
the \ImNScreen{} unattributed crossings lying in a catalogued block &
\ImNNoParent{} have no catalogued transition in the image interval and
\ImNParentAbsent{} have one the same block observed with nothing bright
enough to leak; for \ImNLeft{} the parent frequency was not observed &
-- \\
\hline
\end{tabular}
\end{table}



\emph{One execution block fails these tests, and it fails them as a block
and not as a carrier.} Its \DqBlockNCross{} crossings toward
\LgVerRejectStar{} are the \DqBlockNCross{} highest statistics in the survey
(\S\ref{sec:bdblock}), and three readings say interference does not account
for them: the visibility fit returns a large imaginary part, placing the
emission away from the phase centre; the block's \BdNSibling{} sibling
executions of the same field, array and spectral setup produce no crossing at
all; and of the further windows covering each of its frequencies toward other
stars, \DqOtherSame{} carries a crossing in the same channel.
\emph{None of this substitutes for the ON/OFF cadence a dedicated search
uses}: allocation tables cannot exclude harmonics, intermodulation products,
unconventional emitters or near-field sources, and archival data without a
simultaneous off-source strategy cannot do better. These are statements about
the population and not a clearance of any one crossing.

%% "The null result rests on internal validation alone" lives once, in the
%% discussion's lessons subsection, and not here.  fig:appm is not included
%% -- it compared our own two statistics and the paper states only the one it
%% uses -- so figures/appm_power.pdf and the \AppMPowerOld* and
%% \RfiOcc{Prior,Dir,Stale}* families will be retired.
%%
%% ★ r16-mask: the cross-target screen of the fourth test is now scoped to
%% the ADOPTED attribution (round 660, maskrfi_v660.py), which makes it two
%% events and not one.  round 49's \RfiEventLead, \RfiFreqThem,
%% \RfiSameChanClause, \RfiNotTestableClause, \RfiOtherWin,
%% \RfiOtherStarsRange and \RfiOtherBlocksRange are therefore unreferenced
%% and will be retired; rfi_v399.py still reproduces them and its grammar
%% selftest still runs, so nothing breaks if the scope is ever restored.
